\section{引言}

Nicolai 映射把相互作用的超对称理论
变换为无相互作用的理论~\cite{NicolaiMap}。
人们认为超对称是这类场变换得以存在的关键，
原因在于其雅可比行列式恰好被费米子配分函数所抵消。

在格点规范理论中，
一个自然的问题是：
是否存在把理论映射到其强耦合极限的场变换。
特别地，在没有物质场的情况下，
人们寻找的是对泛函积分中规范场 $U$ 的代换
\begin{equation}
  U=\trans(V)
\end{equation}
使其雅可比行列式抵消规范场作用量。
与 Nicolai 映射类似，这类变换把理论映射到一个可解的理论，
但并不需要超对称，且雅可比行列式所起的作用也不相同。

在有限格点上，若规范群是紧致且连通的，则由紧致流形上体积形式的一般定理
（例如参见文献~\cite{MoserZehnder}，定理~1.26）
可保证这类平凡化映射的存在性。
人们可能倾向于认为这些变换过于复杂而没有实用价值。
然而，正如本文稍后所解释的，
可以通过对场空间中的流作积分来逐步构建平凡化映射，
其生成元满足某些偏微分方程。
这些方程相当易于处理，
并且在纯规范理论中在一定程度上可以解析求解。
由此可以设想的平凡化映射的一个应用，
便是加速格点 QCD 模拟。

\begin{figure}[htbp]
  \centering
  \includegraphics[width=6.0cm]{images/cycle.pdf}
  \caption{%
本文提出的格点 QCD 模拟算法
按照图中箭头所示分三步更新规范场 $U$，其中
HMC 步骤使用的哈密顿函数
包含标准动能项以及场变换 $\mathcal{F}$ 的雅可比行列式
（参见 2.4 小节）。
}
  \label{fig:cycle}
\end{figure}

现有模拟算法的效率不可预测，
这一直是数值格点 QCD 的一个弱点。
不久之前，Del Debbio、Panagopoulos 和 Vicari~\cite{DelDebbioTauQ}
对 $\Group$ 规范理论的经验研究表明，
与规范场拓扑荷相关的可观测量
的自关联时间往往很大，
并且看起来随格距的倒数呈指数增长。
此外，
Schaefer、Sommer 和 Virotta~\cite{StefanConference}
最近发现，就这方面而言，
当模拟中包含海夸克时，情况基本没有改变。

通过将近似平凡化映射
与 HMC 模拟算法~\cite{HMC} 相结合，
小格距下模拟的急剧减速或许可以克服
（见图~1）。
由于该变换使理论更接近强耦合极限——已知 HMC 算法在该极限下效率很高——
自关联时间一般预期会以这种方式缩短。
显然，要让组合算法在实际中奏效，
必须找到相当简单且易于编程实现的近似平凡化映射。
因此，本文的目标之一就是为这一问题提供一个解决方案。
